\begin{ExerciseBlock}
\ExerciseHeading{01}{Lichtbrechung am Materialübergang}{geometrical_and_wave_optics_water_air_refraction}
Ein Lichtstrahl trifft aus Wasser ($n_\mathrm W=1{,}33$) auf die Grenzfläche zur Luft ($n_\mathrm L\approx 1$). Alle Winkel werden zum Lot gemessen.
\begin{center}
\CachedTikzFigure{content/resources/chapters/geometrical_and_wave_optics/tikz/water_air_interface.tex}
\end{center}
\begin{enumerate}
\item Berechnen Sie die Ausbreitungsgeschwindigkeit $c_\mathrm W$ im Wasser.
\item Der Einfallswinkel beträgt $\vartheta_\mathrm W=\ang{20}$. Berechnen Sie den Reflexionswinkel $\vartheta_\mathrm R$ und den Austrittswinkel $\vartheta_\mathrm L$.
\item Für welche Einfallswinkel $\vartheta_\mathrm W$ tritt Totalreflexion auf?
\item Kann ein Lichtstrahl, der aus der Luft auf die Wasseroberfläche trifft, totalreflektiert werden?
\end{enumerate}
\end{ExerciseBlock}
\UseExerciseSolution{geometrical_and_wave_optics_water_air_refraction}
